\section*{A2.3. Registro de supuestos}
Este anexo reúne los supuestos de formulación y colaboración, identificados como SUP-01 a SUP-05, y las 22 decisiones tomadas por el proponente para formular la solución, identificadas como DEC-01 a DEC-22. Estas decisiones forman parte del registro de supuestos y conservan sus identificadores para mantener la trazabilidad.

Cada supuesto o decisión explicita el criterio asumido, su fundamento, el impacto si resulta equivocado y la instancia prevista de validación. Su inclusión no acredita aceptación del cliente ni validación ejecutada.

La responsabilidad de formular las decisiones técnicas corresponde al proponente. La validación con el cliente confirma condiciones operacionales y atribuciones de negocio; no sustituye el diseño. La decisión DEC-17 corresponde a las medidas de cobranza sobre morosos contempladas en el documento del Caso 07.


\section*{A2.3.1 Supuestos de formulación y colaboración}
Las fichas siguientes identifican los supuestos, su fundamento, impacto y validación prevista.

\subsection*{SUP-01 — Disponibilidad de representantes funcionales}
La ficha establece el criterio de trabajo y las condiciones necesarias para su validación.
La Tabla~\ref{tab:a2-3-01} detalla los campos y las condiciones de validación de esta ficha.
\begin{longtable}{|L{.25\linewidth}|L{\dimexpr.75\linewidth-4\tabcolsep-3\arrayrulewidth\relax}|}
\caption{SUP-01 — Disponibilidad de representantes funcionales}\label{tab:a2-3-01}\\
\hline
\EncabezadoTabla{Campo} & \EncabezadoTabla{Contenido} \\ \hline
\endfirsthead
\multicolumn{2}{l}{\small\textit{SUP-01 — Disponibilidad de representantes funcionales — continuación}}\\
\hline
\EncabezadoTabla{Campo} & \EncabezadoTabla{Contenido} \\ \hline
\endhead
\multicolumn{2}{r}{\small\textit{}}\\
\endfoot
\endlastfoot
\rowcolor{RFclaro}
Supuesto & El cliente facilitará participación acotada de representantes de operaciones, administración, escuela y varadero para revisar procesos y validar resultados, sin suponer dedicación exclusiva. \\ \hline
Fundamento & La validación necesita conocimiento operativo y la compañía tiene dotación limitada. \\ \hline
\rowcolor{RFclaro}
Impacto si no se cumple & Decisiones tardías, reprogramación de talleres y riesgo de aceptación; el proponente deberá planificar alternativas y escalamiento. \\ \hline
Instancia de validación & Al inicio de planificación y antes de cada validación: acordar representantes, suplentes y disponibilidad por actividad. \\ \hline
\rowcolor{RFclaro}
Responsable propuesto & Responsable del proyecto y contraparte del cliente. \\ \hline
Fuente y referencia & SD2 sec. 2.5.2; Evidencia por incorporar: RACI y agenda acordada en SD6/SD7. \\ \hline
\rowcolor{RFclaro}
Estado & Validación documental no acreditada; registrar responsable, resultado y tratamiento del desvío. \\ \hline
\end{longtable}
\FuenteTabla{Registro de Synaptix; fuentes y vínculos identificados en las filas.}

El impacto y la instancia de validación permiten priorizar el control de este antecedente.
\subsection*{SUP-02 — Acceso autorizado a fuentes}
La ficha establece el criterio de trabajo y las condiciones necesarias para su validación.
La Tabla~\ref{tab:a2-3-02} detalla los campos y las condiciones de validación de esta ficha.
\begin{longtable}{|L{.25\linewidth}|L{\dimexpr.75\linewidth-4\tabcolsep-3\arrayrulewidth\relax}|}
\caption{SUP-02 — Acceso autorizado a fuentes}\label{tab:a2-3-02}\\
\hline
\EncabezadoTabla{Campo} & \EncabezadoTabla{Contenido} \\ \hline
\endfirsthead
\multicolumn{2}{l}{\small\textit{SUP-02 — Acceso autorizado a fuentes — continuación}}\\
\hline
\EncabezadoTabla{Campo} & \EncabezadoTabla{Contenido} \\ \hline
\endhead
\multicolumn{2}{r}{\small\textit{}}\\
\endfoot
\endlastfoot
\rowcolor{RFclaro}
Supuesto & El cliente permitirá el acceso autorizado a sistemas, archivos y documentación necesarios, dentro de sus atribuciones; no se presume calidad, integridad, documentación ni disponibilidad de APIs. \\ \hline
Fundamento & La migración e integración requieren examinar datos reales y permisos. \\ \hline
\rowcolor{RFclaro}
Impacto si no se cumple & Mayor esfuerzo de extracción, saneamiento y conciliación o necesidad de replanificar; no justifica omitir datos obligatorios. \\ \hline
Instancia de validación & Antes del perfilado e integración: verificar permisos, inventario, extracción y disponibilidad del único conocedor del legado. \\ \hline
\rowcolor{RFclaro}
Responsable propuesto & Responsable de datos y contraparte propietaria de la información. \\ \hline
Fuente y referencia & SD2 sec. 2.5.2; C07 RT-05.15; Evidencia por incorporar: inventario y autorizaciones vigentes en SD5. \\ \hline
\rowcolor{RFclaro}
Estado & Validación documental no acreditada; registrar responsable, resultado y tratamiento del desvío. \\ \hline
\end{longtable}
\FuenteTabla{Registro de Synaptix; fuentes y vínculos identificados en las filas.}

El impacto y la instancia de validación permiten priorizar el control de este antecedente.
\subsection*{SUP-03 — Volumetría como línea base}
La ficha establece el criterio de trabajo y las condiciones necesarias para su validación.
La Tabla~\ref{tab:a2-3-03} detalla los campos y las condiciones de validación de esta ficha.
\begin{longtable}{|L{.25\linewidth}|L{\dimexpr.75\linewidth-4\tabcolsep-3\arrayrulewidth\relax}|}
\caption{SUP-03 — Volumetría como línea base}\label{tab:a2-3-03}\\
\hline
\EncabezadoTabla{Campo} & \EncabezadoTabla{Contenido} \\ \hline
\endfirsthead
\multicolumn{2}{l}{\small\textit{SUP-03 — Volumetría como línea base — continuación}}\\
\hline
\EncabezadoTabla{Campo} & \EncabezadoTabla{Contenido} \\ \hline
\endhead
\multicolumn{2}{r}{\small\textit{}}\\
\endfoot
\endlastfoot
\rowcolor{RFclaro}
Supuesto & Los valores del caso constituyen la línea base de dimensionamiento; personas, embarcaciones y movimientos no equivalen automáticamente a usuarios concurrentes, mensajes o transacciones. \\ \hline
Fundamento & El peak y las dependencias requieren derivación explícita de carga. \\ \hline
\rowcolor{RFclaro}
Impacto si no se cumple & Subdimensionamiento o sobredimensionamiento si la conversión a carga es incorrecta. \\ \hline
Instancia de validación & Durante formulación de arquitectura y antes de pruebas: declarar fórmulas, escenarios y crecimiento, y contrastar con pruebas. \\ \hline
\rowcolor{RFclaro}
Responsable propuesto & Responsable de arquitectura y operaciones. \\ \hline
Fuente y referencia & SD2 sec. 2.5.2; C07 sec. 14.1–14.2 y RT-09.02; Evidencia por incorporar: dimensionamiento SD4 y plan de carga SD9. \\ \hline
\rowcolor{RFclaro}
Estado & Validación documental no acreditada; registrar responsable, resultado y tratamiento del desvío. \\ \hline
\end{longtable}
\FuenteTabla{Registro de Synaptix; fuentes y vínculos identificados en las filas.}

El impacto y la instancia de validación permiten priorizar el control de este antecedente.
\subsection*{SUP-04 — Compras y obras a cargo del cliente}
La ficha establece el criterio de trabajo y las condiciones necesarias para su validación.
La Tabla~\ref{tab:a2-3-04} detalla los campos y las condiciones de validación de esta ficha.
\begin{longtable}{|L{.25\linewidth}|L{\dimexpr.75\linewidth-4\tabcolsep-3\arrayrulewidth\relax}|}
\caption{SUP-04 — Compras y obras a cargo del cliente}\label{tab:a2-3-04}\\
\hline
\EncabezadoTabla{Campo} & \EncabezadoTabla{Contenido} \\ \hline
\endfirsthead
\multicolumn{2}{l}{\small\textit{SUP-04 — Compras y obras a cargo del cliente — continuación}}\\
\hline
\EncabezadoTabla{Campo} & \EncabezadoTabla{Contenido} \\ \hline
\endhead
\multicolumn{2}{r}{\small\textit{}}\\
\endfoot
\endlastfoot
\rowcolor{RFclaro}
Supuesto & El cliente ejecutará compras de hardware y obras excluidas de la prestación de construcción, a partir de especificaciones, cantidades, costos y coordinación provistos por el proponente; no se presume infraestructura lista al inicio. \\ \hline
Fundamento & La obligación de especificar e integrar subsiste aunque compre o construya el cliente. \\ \hline
\rowcolor{RFclaro}
Impacto si no se cumple & Retraso de pruebas o despliegue por disponibilidad de equipos, energía o canalizaciones; requiere dependencias y contingencias. \\ \hline
Instancia de validación & Antes de fijar ventanas de instalación: acordar hitos de adquisición, recepción y habilitación con responsables. \\ \hline
\rowcolor{RFclaro}
Responsable propuesto & Responsable de infraestructura y contraparte de compras/obras del cliente. \\ \hline
Fuente y referencia & SD2 sec. 2.5.2; C07 cap. 11; Evidencia por incorporar: SD4/T-11 y dependencias del cronograma SD7. \\ \hline
\rowcolor{RFclaro}
Estado & Validación documental no acreditada; registrar responsable, resultado y tratamiento del desvío. \\ \hline
\end{longtable}
\FuenteTabla{Registro de Synaptix; fuentes y vínculos identificados en las filas.}

El impacto y la instancia de validación permiten priorizar el control de este antecedente.
\subsection*{SUP-05 — Aprobación de decisiones de negocio}
La ficha establece el criterio de trabajo y las condiciones necesarias para su validación.
La Tabla~\ref{tab:a2-3-05} detalla los campos y las condiciones de validación de esta ficha.
\begin{longtable}{|L{.25\linewidth}|L{\dimexpr.75\linewidth-4\tabcolsep-3\arrayrulewidth\relax}|}
\caption{SUP-05 — Aprobación de decisiones de negocio}\label{tab:a2-3-05}\\
\hline
\EncabezadoTabla{Campo} & \EncabezadoTabla{Contenido} \\ \hline
\endfirsthead
\multicolumn{2}{l}{\small\textit{SUP-05 — Aprobación de decisiones de negocio — continuación}}\\
\hline
\EncabezadoTabla{Campo} & \EncabezadoTabla{Contenido} \\ \hline
\endhead
\multicolumn{2}{r}{\small\textit{}}\\
\endfoot
\endlastfoot
\rowcolor{RFclaro}
Supuesto & El cliente designará las autoridades competentes para aprobar criterios comerciales y operativos; el proponente formulará y justificará sus propuestas, sin trasladar al cliente las 22 decisiones de diseño ni atribuirse facultades societarias. \\ \hline
Fundamento & El diseño requiere criterios concretos y los actos del cliente necesitan atribuciones. \\ \hline
\rowcolor{RFclaro}
Impacto si no se cumple & Reglas sin legitimidad, cambios reiterados o uso de medidas no autorizadas. \\ \hline
Instancia de validación & Durante formulación y antes de parametrizar reglas: identificar autoridad y registrar decisión, alcance y versión. \\ \hline
\rowcolor{RFclaro}
Responsable propuesto & Responsable del proyecto y autoridad de negocio competente. \\ \hline
Fuente y referencia & SD2 sec. 2.5.2; C07 sec. 16.1; Evidencia por incorporar: matriz de atribuciones y actas de aprobación. \\ \hline
\rowcolor{RFclaro}
Estado & Validación documental no acreditada; registrar responsable, resultado y tratamiento del desvío. \\ \hline
\end{longtable}
\FuenteTabla{Registro de Synaptix; fuentes y vínculos identificados en las filas.}

El impacto y la instancia de validación permiten priorizar el control de este antecedente.


\section*{A2.3.2 Supuestos de decisiones del proponente}
Las fichas siguientes registran las decisiones de formulación y sus condiciones de validación.

\subsection*{DEC-01 — Estado de ocupación y salida}
La ficha establece el criterio de trabajo y las condiciones necesarias para su validación.
La Tabla~\ref{tab:a2-3-06} detalla los campos y las condiciones de validación de esta ficha.
\begin{longtable}{|L{.25\linewidth}|L{\dimexpr.75\linewidth-4\tabcolsep-3\arrayrulewidth\relax}|}
\caption{DEC-01 — Estado de ocupación y salida}\label{tab:a2-3-06}\\
\hline
\EncabezadoTabla{Campo} & \EncabezadoTabla{Contenido} \\ \hline
\endfirsthead
\multicolumn{2}{l}{\small\textit{DEC-01 — Estado de ocupación y salida — continuación}}\\
\hline
\EncabezadoTabla{Campo} & \EncabezadoTabla{Contenido} \\ \hline
\endhead
\multicolumn{2}{r}{\small\textit{}}\\
\endfoot
\endlastfoot
\rowcolor{RFclaro}
Fuente de la decisión & C07, sec. 16.1, decisión 1. \\ \hline
Requerimientos relacionados & RF-NAV-01, RF-NAV-04, RF-AMA-03 \\ \hline
\rowcolor{RFclaro}
Criterio formulado para revisión & Combinar confirmación identificada del marinero o guardia y declaraciones del armador; conservar discrepancias y antigüedad de evidencia. La elección concreta del mecanismo de detección requerido por RT-17.06 permanece abierta hasta comparar fiabilidad, cobertura y costo; no declarar resuelto el diseño completo. \\ \hline
Fundamento & Evita equiparar ausencia de registro con ausencia física y mantiene la radio. \\ \hline
\rowcolor{RFclaro}
Impacto si el criterio resulta incorrecto & Una ocupación falsa altera asignación, avisos y retorno seguro. \\ \hline
Instancia y responsables propuestos de validación & Operaciones marítimas y responsable de arquitectura: prueba de campo y comparación de alternativas antes de cerrar arquitectura. \\ \hline
\rowcolor{RFclaro}
Estado de formulación & Parcial: mecanismo de detección sin evidencia suficiente. No se presenta como acuerdo del equipo. \\ \hline
Referencia faltante & Evidencia por incorporar: SD3/SD4: decisión del mecanismo, precisión y manejo de discrepancias; T-11 si incorpora sensores. \\ \hline
\end{longtable}
\FuenteTabla{Registro de Synaptix; fuentes y vínculos identificados en las filas.}

El impacto y la instancia de validación permiten priorizar el control de este antecedente.
\subsection*{DEC-02 — Legado y migración}
La ficha establece el criterio de trabajo y las condiciones necesarias para su validación.
La Tabla~\ref{tab:a2-3-07} detalla los campos y las condiciones de validación de esta ficha.
\begin{longtable}{|L{.25\linewidth}|L{\dimexpr.75\linewidth-4\tabcolsep-3\arrayrulewidth\relax}|}
\caption{DEC-02 — Legado y migración}\label{tab:a2-3-07}\\
\hline
\EncabezadoTabla{Campo} & \EncabezadoTabla{Contenido} \\ \hline
\endfirsthead
\multicolumn{2}{l}{\small\textit{DEC-02 — Legado y migración — continuación}}\\
\hline
\EncabezadoTabla{Campo} & \EncabezadoTabla{Contenido} \\ \hline
\endhead
\multicolumn{2}{r}{\small\textit{}}\\
\endfoot
\endlastfoot
\rowcolor{RFclaro}
Fuente de la decisión & C07, sec. 16.1, decisión 2. \\ \hline
Requerimientos relacionados & RF-ADM-06, RF-ADM-07, RNF-MIG-01 \\ \hline
\rowcolor{RFclaro}
Criterio formulado para revisión & Migrar íntegramente el alcance obligatorio de RT-05.15, preservar otros antecedentes exigibles en consulta protegida y ejecutar dos ensayos, conciliación y reversión; no conservar el software obsoleto como única vía de consulta. \\ \hline
Fundamento & Distingue los doce años de historia del alcance mínimo de migración y de la retención. \\ \hline
\rowcolor{RFclaro}
Impacto si el criterio resulta incorrecto & Pérdida de saldos o evidencia de custodia impide cobro y defensa documental. \\ \hline
Instancia y responsables propuestos de validación & Administración, conocedor del legado y responsable de datos: perfilado y conciliación antes de cada corte. \\ \hline
\rowcolor{RFclaro}
Estado de formulación & Criterio parcialmente documentado; desarrollo pendiente. No se presenta como acuerdo del equipo. \\ \hline
Referencia faltante & Evidencia por incorporar: SD5, plan de migración, inventario de fuentes y acta de conciliación. \\ \hline
\end{longtable}
\FuenteTabla{Registro de Synaptix; fuentes y vínculos identificados en las filas.}

El impacto y la instancia de validación permiten priorizar el control de este antecedente.
\subsection*{DEC-03 — Medición por amarra}
La ficha establece el criterio de trabajo y las condiciones necesarias para su validación.
La Tabla~\ref{tab:a2-3-08} detalla los campos y las condiciones de validación de esta ficha.
\begin{longtable}{|L{.25\linewidth}|L{\dimexpr.75\linewidth-4\tabcolsep-3\arrayrulewidth\relax}|}
\caption{DEC-03 — Medición por amarra}\label{tab:a2-3-08}\\
\hline
\EncabezadoTabla{Campo} & \EncabezadoTabla{Contenido} \\ \hline
\endfirsthead
\multicolumn{2}{l}{\small\textit{DEC-03 — Medición por amarra — continuación}}\\
\hline
\EncabezadoTabla{Campo} & \EncabezadoTabla{Contenido} \\ \hline
\endhead
\multicolumn{2}{r}{\small\textit{}}\\
\endfoot
\endlastfoot
\rowcolor{RFclaro}
Fuente de la decisión & C07, sec. 16.1, decisión 3. \\ \hline
Requerimientos relacionados & RF-CON-01 a RF-CON-06 \\ \hline
\rowcolor{RFclaro}
Criterio formulado para revisión & Proponer medición atribuible a cada amarra, con asociación temporal a contrato y una misma serie para cobro y ambiente. No consolidar como decisión aceptada el reparto de una torreta: requiere fórmula, precisión, tratamiento de cambios de ocupante y respaldo del marco de traspaso de costos. \\ \hline
Fundamento & Evita cobrar consumos ajenos y construir indicadores con una serie diferente. \\ \hline
\rowcolor{RFclaro}
Impacto si el criterio resulta incorrecto & Cobros discutibles, errores de atribución y falta de serie ambiental. \\ \hline
Instancia y responsables propuestos de validación & Responsable de integración, administración y especialista competente: piloto en pantalán y revisión del marco aplicable antes de especificar compra. \\ \hline
\rowcolor{RFclaro}
Estado de formulación & Abierta la selección física y el respaldo del cobro. No se presenta como acuerdo del equipo. \\ \hline
Referencia faltante & Evidencia por incorporar: SD4/T-11: esquema y cantidad de medidores; SD5: datos; SD11: costos; investigación sectorial. \\ \hline
\end{longtable}
\FuenteTabla{Registro de Synaptix; fuentes y vínculos identificados en las filas.}

El impacto y la instancia de validación permiten priorizar el control de este antecedente.
\subsection*{DEC-04 — Alcance del plan voluntario}
La ficha establece el criterio de trabajo y las condiciones necesarias para su validación.
La Tabla~\ref{tab:a2-3-09} detalla los campos y las condiciones de validación de esta ficha.
\begin{longtable}{|L{.25\linewidth}|L{\dimexpr.75\linewidth-4\tabcolsep-3\arrayrulewidth\relax}|}
\caption{DEC-04 — Alcance del plan voluntario}\label{tab:a2-3-09}\\
\hline
\EncabezadoTabla{Campo} & \EncabezadoTabla{Contenido} \\ \hline
\endfirsthead
\multicolumn{2}{l}{\small\textit{DEC-04 — Alcance del plan voluntario — continuación}}\\
\hline
\EncabezadoTabla{Campo} & \EncabezadoTabla{Contenido} \\ \hline
\endhead
\multicolumn{2}{r}{\small\textit{}}\\
\endfoot
\endlastfoot
\rowcolor{RFclaro}
Fuente de la decisión & C07, sec. 16.1, decisión 4. \\ \hline
Requerimientos relacionados & RF-NAV-06 a RF-NAV-09, RF-AUT-08 \\ \hline
\rowcolor{RFclaro}
Criterio formulado para revisión & Proponer campos de nave, patrón, personas a bordo, salida, destino o recorrido declarado, retorno previsto y contactos; permitir actualización y cierre identificados. Alertar dentro de 15 minutos desde vencimiento sin gracia adicional; derivar a guardia y aplicar protocolo documentado. La marina recibe antecedentes y comunica, sin prometer vigilancia permanente ni búsqueda y salvamento. \\ \hline
Fundamento & Define el compromiso y evita una expectativa de rescate automático. \\ \hline
\rowcolor{RFclaro}
Impacto si el criterio resulta incorrecto & Demora de aviso, falsa tranquilidad o obligación operacional no financiada. \\ \hline
Instancia y responsables propuestos de validación & Operaciones y responsable de servicio: escenario de vencimiento, mensaje tardío y falso cierre antes de validar el flujo. \\ \hline
\rowcolor{RFclaro}
Estado de formulación & Propuesta completada para revisión; completar su justificación. No se presenta como acuerdo del equipo. \\ \hline
Referencia faltante & Evidencia por incorporar: SD3: alcance; SD4: estados y canal de baja capacidad; protocolo de escalamiento y destinatarios. \\ \hline
\end{longtable}
\FuenteTabla{Registro de Synaptix; fuentes y vínculos identificados en las filas.}

El impacto y la instancia de validación permiten priorizar el control de este antecedente.
\subsection*{DEC-05 — Confirmación del regreso}
La ficha establece el criterio de trabajo y las condiciones necesarias para su validación.
La Tabla~\ref{tab:a2-3-10} detalla los campos y las condiciones de validación de esta ficha.
\begin{longtable}{|L{.25\linewidth}|L{\dimexpr.75\linewidth-4\tabcolsep-3\arrayrulewidth\relax}|}
\caption{DEC-05 — Confirmación del regreso}\label{tab:a2-3-10}\\
\hline
\EncabezadoTabla{Campo} & \EncabezadoTabla{Contenido} \\ \hline
\endfirsthead
\multicolumn{2}{l}{\small\textit{DEC-05 — Confirmación del regreso — continuación}}\\
\hline
\EncabezadoTabla{Campo} & \EncabezadoTabla{Contenido} \\ \hline
\endhead
\multicolumn{2}{r}{\small\textit{}}\\
\endfoot
\endlastfoot
\rowcolor{RFclaro}
Fuente de la decisión & C07, sec. 16.1, decisión 5. \\ \hline
Requerimientos relacionados & RF-NAV-04, RF-NAV-05 \\ \hline
\rowcolor{RFclaro}
Criterio formulado para revisión & El marinero identificado confirma la embarcación amarrada en condición segura; la guardia puede registrar la observación comunicada por radio, identificando informante y registrador. El cierre declarado por el armador se conserva como evidencia y se concilia con la confirmación operacional; una discrepancia no se borra automáticamente. \\ \hline
Fundamento & Permite cierre sin dispositivo a bordo y separa autor de la evidencia de quien digita. \\ \hline
\rowcolor{RFclaro}
Impacto si el criterio resulta incorrecto & Cierre falso o persistencia de alertas pese al regreso. \\ \hline
Instancia y responsables propuestos de validación & Jefatura de operaciones y marineros: prueba de regreso, amarra distinta y reporte por tercero en momento seguro. \\ \hline
\rowcolor{RFclaro}
Estado de formulación & Criterio formulado; pendiente ratificación. No se presenta como acuerdo del equipo. \\ \hline
Referencia faltante & Evidencia por incorporar: SD3/SD4: autoridad de cierre, catálogo de evidencias y manejo de discrepancias. \\ \hline
\end{longtable}
\FuenteTabla{Registro de Synaptix; fuentes y vínculos identificados en las filas.}

El impacto y la instancia de validación permiten priorizar el control de este antecedente.
\subsection*{DEC-06 — Alumno, bote y monitor}
La ficha establece el criterio de trabajo y las condiciones necesarias para su validación.
La Tabla~\ref{tab:a2-3-11} detalla los campos y las condiciones de validación de esta ficha.
\begin{longtable}{|L{.25\linewidth}|L{\dimexpr.75\linewidth-4\tabcolsep-3\arrayrulewidth\relax}|}
\caption{DEC-06 — Alumno, bote y monitor}\label{tab:a2-3-11}\\
\hline
\EncabezadoTabla{Campo} & \EncabezadoTabla{Contenido} \\ \hline
\endfirsthead
\multicolumn{2}{l}{\small\textit{DEC-06 — Alumno, bote y monitor — continuación}}\\
\hline
\EncabezadoTabla{Campo} & \EncabezadoTabla{Contenido} \\ \hline
\endhead
\multicolumn{2}{r}{\small\textit{}}\\
\endfoot
\endlastfoot
\rowcolor{RFclaro}
Fuente de la decisión & C07, sec. 16.1, decisión 6. \\ \hline
Requerimientos relacionados & RF-ESC-03 a RF-ESC-06, RF-ESC-09 \\ \hline
\rowcolor{RFclaro}
Criterio formulado para revisión & Preparar nómina y asignación antes de salir; confirmar grupalmente en momento seguro, registrar cambios de bote o monitor y verificar el regreso de cada alumno antes del cierre. Sin cámaras de alumnos ni dispositivos en embarcaciones menores; sin exigir interacción durante una maniobra. \\ \hline
Fundamento & Produce evidencia humana identificada dentro del plazo de 60 segundos para hasta 18 botes. \\ \hline
\rowcolor{RFclaro}
Impacto si el criterio resulta incorrecto & No saber quién permanece en el agua y confundir asistencia prevista con efectiva. \\ \hline
Instancia y responsables propuestos de validación & Coordinación y monitores: ensayo realista con lluvia, ausencias y cambios antes de recepción del módulo. \\ \hline
\rowcolor{RFclaro}
Estado de formulación & Criterio documentado, detalle por validar. No se presenta como acuerdo del equipo. \\ \hline
Referencia faltante & Evidencia por incorporar: SD3/SD4: estados y pantalla; SD9: prueba de clase completa. \\ \hline
\end{longtable}
\FuenteTabla{Registro de Synaptix; fuentes y vínculos identificados en las filas.}

El impacto y la instancia de validación permiten priorizar el control de este antecedente.
\subsection*{DEC-07 — Retiro anticipado}
La ficha establece el criterio de trabajo y las condiciones necesarias para su validación.
La Tabla~\ref{tab:a2-3-12} detalla los campos y las condiciones de validación de esta ficha.
\begin{longtable}{|L{.25\linewidth}|L{\dimexpr.75\linewidth-4\tabcolsep-3\arrayrulewidth\relax}|}
\caption{DEC-07 — Retiro anticipado}\label{tab:a2-3-12}\\
\hline
\EncabezadoTabla{Campo} & \EncabezadoTabla{Contenido} \\ \hline
\endfirsthead
\multicolumn{2}{l}{\small\textit{DEC-07 — Retiro anticipado — continuación}}\\
\hline
\EncabezadoTabla{Campo} & \EncabezadoTabla{Contenido} \\ \hline
\endhead
\multicolumn{2}{r}{\small\textit{}}\\
\endfoot
\endlastfoot
\rowcolor{RFclaro}
Fuente de la decisión & C07, sec. 16.1, decisión 7. \\ \hline
Requerimientos relacionados & RF-ESC-07, RF-ESC-08, RF-ESC-14 \\ \hline
\rowcolor{RFclaro}
Criterio formulado para revisión & Usar un único expediente de asistencia con retiro solicitado, verificado y confirmado; portería solicita y comprueba identidad y autorización, escuela verifica situación y entrega segura, y ambos confirman el mismo evento. Si falla la comunicación digital, confirmar por radio y registrar la contingencia; se identifica quién recibe la confirmación y quién la registra, se conserva el vínculo con el evento para la reconciliación posterior. \\ \hline
Fundamento & Resuelve la divergencia de dos registros; una alerta sola no prueba entrega. \\ \hline
\rowcolor{RFclaro}
Impacto si el criterio resulta incorrecto & Alumno retirado permanece como presente o alumno en agua se da por entregado. \\ \hline
Instancia y responsables propuestos de validación & Portería, escuela y responsable de continuidad: ensayo de retiro con clase en agua, WAN caída y terminal aislado. \\ \hline
\rowcolor{RFclaro}
Estado de formulación & Propuesta de corrección; validación operacional imprescindible. No se presenta como acuerdo del equipo. \\ \hline
Referencia faltante & Evidencia por incorporar: SD3/SD4: protocolo único y consistencia; SD9: prueba de fallo simultáneo. \\ \hline
\end{longtable}
\FuenteTabla{Registro de Synaptix; fuentes y vínculos identificados en las filas.}

El impacto y la instancia de validación permiten priorizar el control de este antecedente.
\subsection*{DEC-08 — Vista de salud mínima}
La ficha establece el criterio de trabajo y las condiciones necesarias para su validación.
La Tabla~\ref{tab:a2-3-13} detalla los campos y las condiciones de validación de esta ficha.
\begin{longtable}{|L{.25\linewidth}|L{\dimexpr.75\linewidth-4\tabcolsep-3\arrayrulewidth\relax}|}
\caption{DEC-08 — Vista de salud mínima}\label{tab:a2-3-13}\\
\hline
\EncabezadoTabla{Campo} & \EncabezadoTabla{Contenido} \\ \hline
\endfirsthead
\multicolumn{2}{l}{\small\textit{DEC-08 — Vista de salud mínima — continuación}}\\
\hline
\EncabezadoTabla{Campo} & \EncabezadoTabla{Contenido} \\ \hline
\endhead
\multicolumn{2}{r}{\small\textit{}}\\
\endfoot
\endlastfoot
\rowcolor{RFclaro}
Fuente de la decisión & C07, sec. 16.1, decisión 8. \\ \hline
Requerimientos relacionados & RF-ESC-11, RF-ESC-12, RNF-SEG-04 \\ \hline
\rowcolor{RFclaro}
Criterio formulado para revisión & Proponer que el monitor sólo vea alertas y pautas de emergencia necesarias y autorizadas para los alumnos a su cargo, junto con contacto de emergencia. La coordinación designada por el cliente aprueba visibilidad y vigencia con respaldo profesional pertinente y autorización aplicable; el monitor no accede al expediente íntegro ni se inventan pautas clínicas. \\ \hline
Fundamento & Minimiza exposición sin impedir acceso útil en emergencia. \\ \hline
\rowcolor{RFclaro}
Impacto si el criterio resulta incorrecto & Divulgación indebida o falta de información relevante para asistencia. \\ \hline
Instancia y responsables propuestos de validación & Coordinación, responsable de protección de datos y contraparte competente: revisión de campos y prueba de permisos antes de carga real. \\ \hline
\rowcolor{RFclaro}
Estado de formulación & Parcial: campos y autoridad necesitan referencia. No se presenta como acuerdo del equipo. \\ \hline
Referencia faltante & Evidencia por incorporar: SD5: clasificación/campos; SD4: permisos; autorización y fundamento de tratamiento. \\ \hline
\end{longtable}
\FuenteTabla{Registro de Synaptix; fuentes y vínculos identificados en las filas.}

El impacto y la instancia de validación permiten priorizar el control de este antecedente.
\subsection*{DEC-09 — Compatibilidad y excepciones}
La ficha establece el criterio de trabajo y las condiciones necesarias para su validación.
La Tabla~\ref{tab:a2-3-14} detalla los campos y las condiciones de validación de esta ficha.
\begin{longtable}{|L{.25\linewidth}|L{\dimexpr.75\linewidth-4\tabcolsep-3\arrayrulewidth\relax}|}
\caption{DEC-09 — Compatibilidad y excepciones}\label{tab:a2-3-14}\\
\hline
\EncabezadoTabla{Campo} & \EncabezadoTabla{Contenido} \\ \hline
\endfirsthead
\multicolumn{2}{l}{\small\textit{DEC-09 — Compatibilidad y excepciones — continuación}}\\
\hline
\EncabezadoTabla{Campo} & \EncabezadoTabla{Contenido} \\ \hline
\endhead
\multicolumn{2}{r}{\small\textit{}}\\
\endfoot
\endlastfoot
\rowcolor{RFclaro}
Fuente de la decisión & C07, sec. 16.1, decisión 9. \\ \hline
Requerimientos relacionados & RF-AMA-03 a RF-AMA-05 \\ \hline
\rowcolor{RFclaro}
Criterio formulado para revisión & Evaluar eslora, manga, calado, tipo y condiciones de maniobra/exposición; bloquear incompatibilidades de seguridad. Someter excepciones comerciales admisibles a responsable de operaciones y segundo perfil habilitado, con razón y evidencia; registrar la matriz de límites antes de parametrizar. \\ \hline
Fundamento & La aprobación jerárquica no vuelve segura una amarra incompatible. \\ \hline
\rowcolor{RFclaro}
Impacto si el criterio resulta incorrecto & Daño a embarcaciones o asignaciones arbitrarias. \\ \hline
Instancia y responsables propuestos de validación & Jefatura de operaciones: validar matriz con casos límite y trazabilidad de excepción. \\ \hline
\rowcolor{RFclaro}
Estado de formulación & Criterio precisado; límites faltantes. No se presenta como acuerdo del equipo. \\ \hline
Referencia faltante & Evidencia por incorporar: SD3/SD4: reglas y cargos; SD5: atributos; anexo de reglas futuro. \\ \hline
\end{longtable}
\FuenteTabla{Registro de Synaptix; fuentes y vínculos identificados en las filas.}

El impacto y la instancia de validación permiten priorizar el control de este antecedente.
\subsection*{DEC-10 — Lista de espera}
La ficha establece el criterio de trabajo y las condiciones necesarias para su validación.
La Tabla~\ref{tab:a2-3-15} detalla los campos y las condiciones de validación de esta ficha.
\begin{longtable}{|L{.25\linewidth}|L{\dimexpr.75\linewidth-4\tabcolsep-3\arrayrulewidth\relax}|}
\caption{DEC-10 — Lista de espera}\label{tab:a2-3-15}\\
\hline
\EncabezadoTabla{Campo} & \EncabezadoTabla{Contenido} \\ \hline
\endfirsthead
\multicolumn{2}{l}{\small\textit{DEC-10 — Lista de espera — continuación}}\\
\hline
\EncabezadoTabla{Campo} & \EncabezadoTabla{Contenido} \\ \hline
\endhead
\multicolumn{2}{r}{\small\textit{}}\\
\endfoot
\endlastfoot
\rowcolor{RFclaro}
Fuente de la decisión & C07, sec. 16.1, decisión 10. \\ \hline
Requerimientos relacionados & RF-AMA-06, RF-AMA-07 \\ \hline
\rowcolor{RFclaro}
Criterio formulado para revisión & Proponer orden por fecha de ingreso entre candidatos compatibles, con criterio de desempate transparente y registro de ofertas, aceptación, rechazo y vencimiento. FALTA DEFINICIÓN: plazo de respuesta y desempate final; no inventar una subasta o discrecionalidad comercial. \\ \hline
Fundamento & Hace reproducible el orden sin tratar una salida temporal como vacante definitiva. \\ \hline
\rowcolor{RFclaro}
Impacto si el criterio resulta incorrecto & Reclamos de socios y asignación injustificable. \\ \hline
Instancia y responsables propuestos de validación & Administración y autoridad societaria competente: revisar simulación de cola y aprobar regla antes de uso. \\ \hline
\rowcolor{RFclaro}
Estado de formulación & Criterio propuesto formulado parcialmente; plazo de respuesta y tratamiento del vencimiento pendientes de desarrollo. \\ \hline
Referencia faltante & Evidencia por incorporar: SD3: política; reglas de negocio y acta del órgano competente. \\ \hline
\end{longtable}
\FuenteTabla{Registro de Synaptix; fuentes y vínculos identificados en las filas.}

El impacto y la instancia de validación permiten priorizar el control de este antecedente.
\subsection*{DEC-11 — Visitantes sin prepago}
La ficha establece el criterio de trabajo y las condiciones necesarias para su validación.
La Tabla~\ref{tab:a2-3-16} detalla los campos y las condiciones de validación de esta ficha.
\begin{longtable}{|L{.25\linewidth}|L{\dimexpr.75\linewidth-4\tabcolsep-3\arrayrulewidth\relax}|}
\caption{DEC-11 — Visitantes sin prepago}\label{tab:a2-3-16}\\
\hline
\EncabezadoTabla{Campo} & \EncabezadoTabla{Contenido} \\ \hline
\endfirsthead
\multicolumn{2}{l}{\small\textit{DEC-11 — Visitantes sin prepago — continuación}}\\
\hline
\EncabezadoTabla{Campo} & \EncabezadoTabla{Contenido} \\ \hline
\endhead
\multicolumn{2}{r}{\small\textit{}}\\
\endfoot
\endlastfoot
\rowcolor{RFclaro}
Fuente de la decisión & C07, sec. 16.1, decisión 11. \\ \hline
Requerimientos relacionados & RF-AMA-10 a RF-AMA-12, RF-FAC-01 \\ \hline
\rowcolor{RFclaro}
Criterio formulado para revisión & Registrar la llegada y estadía aunque no exista reserva o prepago; asociar responsable de pago, tarifa vigente y evidencia operacional, generar el hecho facturable y derivarlo a contabilidad. Ofrecer pago posterior por canal disponible y seguimiento de pendientes, sin obstaculizar una maniobra segura. \\ \hline
Fundamento & El prepago no cubre a quien llega de noche ni elimina la conciliación. \\ \hline
\rowcolor{RFclaro}
Impacto si el criterio resulta incorrecto & Recaladas no cobradas o exigencias inseguras al arribo. \\ \hline
Instancia y responsables propuestos de validación & Torre y administración: prueba nocturna sin reserva, enlace caído y salida antes de oficina. \\ \hline
\rowcolor{RFclaro}
Estado de formulación & Propuesta complementaria para revisión. No se presenta como acuerdo del equipo. \\ \hline
Referencia faltante & Evidencia por incorporar: SD3/SD4: flujo excepcional; SD5: hecho facturable e integración contable. \\ \hline
\end{longtable}
\FuenteTabla{Registro de Synaptix; fuentes y vínculos identificados en las filas.}

El impacto y la instancia de validación permiten priorizar el control de este antecedente.
\subsection*{DEC-12 — Plan de izaje}
La ficha establece el criterio de trabajo y las condiciones necesarias para su validación.
La Tabla~\ref{tab:a2-3-17} detalla los campos y las condiciones de validación de esta ficha.
\begin{longtable}{|L{.25\linewidth}|L{\dimexpr.75\linewidth-4\tabcolsep-3\arrayrulewidth\relax}|}
\caption{DEC-12 — Plan de izaje}\label{tab:a2-3-17}\\
\hline
\EncabezadoTabla{Campo} & \EncabezadoTabla{Contenido} \\ \hline
\endfirsthead
\multicolumn{2}{l}{\small\textit{DEC-12 — Plan de izaje — continuación}}\\
\hline
\EncabezadoTabla{Campo} & \EncabezadoTabla{Contenido} \\ \hline
\endhead
\multicolumn{2}{r}{\small\textit{}}\\
\endfoot
\endlastfoot
\rowcolor{RFclaro}
Fuente de la decisión & C07, sec. 16.1, decisión 12. \\ \hline
Requerimientos relacionados & RF-VAR-04 a RF-VAR-07 \\ \hline
\rowcolor{RFclaro}
Criterio formulado para revisión & Capturar y versionar antes de la maniobra los antecedentes entregados por fabricante/armador y verificados por personal competente. Operador calificado valida el plan antes de iniciar; encargado controla vigencia y discrepancias. Identificar autor, validador y cambios; no atribuir al software cálculo certificado ni sustituir responsabilidad del ejecutor. \\ \hline
Fundamento & Evita depender de la memoria de una persona y del registro con carga suspendida. \\ \hline
\rowcolor{RFclaro}
Impacto si el criterio resulta incorrecto & Izaje con plan incorrecto, daño y responsabilidad no trazable. \\ \hline
Instancia y responsables propuestos de validación & Encargado de varadero y profesional competente: revisión de ficha y simulación previa a cualquier maniobra real. \\ \hline
\rowcolor{RFclaro}
Estado de formulación & Parcial: responsables y fuente técnica por acreditar. No se presenta como acuerdo del equipo. \\ \hline
Referencia faltante & Evidencia por incorporar: SD3: responsabilidades; SD4: flujo; fuentes técnicas del plan y modalidad de firma. \\ \hline
\end{longtable}
\FuenteTabla{Registro de Synaptix; fuentes y vínculos identificados en las filas.}

El impacto y la instancia de validación permiten priorizar el control de este antecedente.
\subsection*{DEC-13 — Reprogramación por clima}
La ficha establece el criterio de trabajo y las condiciones necesarias para su validación.
La Tabla~\ref{tab:a2-3-18} detalla los campos y las condiciones de validación de esta ficha.
\begin{longtable}{|L{.25\linewidth}|L{\dimexpr.75\linewidth-4\tabcolsep-3\arrayrulewidth\relax}|}
\caption{DEC-13 — Reprogramación por clima}\label{tab:a2-3-18}\\
\hline
\EncabezadoTabla{Campo} & \EncabezadoTabla{Contenido} \\ \hline
\endfirsthead
\multicolumn{2}{l}{\small\textit{DEC-13 — Reprogramación por clima — continuación}}\\
\hline
\EncabezadoTabla{Campo} & \EncabezadoTabla{Contenido} \\ \hline
\endhead
\multicolumn{2}{r}{\small\textit{}}\\
\endfoot
\endlastfoot
\rowcolor{RFclaro}
Fuente de la decisión & C07, sec. 16.1, decisión 13. \\ \hline
Requerimientos relacionados & RF-VAR-01 a RF-VAR-03 \\ \hline
\rowcolor{RFclaro}
Criterio formulado para revisión & Proponer nuevas secuencias considerando marea, viento, disponibilidad del travelift, personal, ubicación de cunas y dependencias; encargado valida antes de confirmar y avisar a armadores y contratistas. Conservar agenda anterior, motivos y avisos pendientes. \\ \hline
Fundamento & Mover fechas sin revisar dependencias puede volver imposible la secuencia. \\ \hline
\rowcolor{RFclaro}
Impacto si el criterio resulta incorrecto & Conflictos de patio, jornadas sobrecargadas y avisos inconsistentes. \\ \hline
Instancia y responsables propuestos de validación & Varadero: simular la cancelación de un día completo de campaña y revalidar dependencias. \\ \hline
\rowcolor{RFclaro}
Estado de formulación & Criterio desarrollado para revisión. No se presenta como acuerdo del equipo. \\ \hline
Referencia faltante & Evidencia por incorporar: SD3/SD4: reglas de replanificación; SD7: reservas de capacidad y calendario. \\ \hline
\end{longtable}
\FuenteTabla{Registro de Synaptix; fuentes y vínculos identificados en las filas.}

El impacto y la instancia de validación permiten priorizar el control de este antecedente.
\subsection*{DEC-14 — Origen de residuos sin app del contratista}
La ficha establece el criterio de trabajo y las condiciones necesarias para su validación.
La Tabla~\ref{tab:a2-3-19} detalla los campos y las condiciones de validación de esta ficha.
\begin{longtable}{|L{.25\linewidth}|L{\dimexpr.75\linewidth-4\tabcolsep-3\arrayrulewidth\relax}|}
\caption{DEC-14 — Origen de residuos sin app del contratista}\label{tab:a2-3-19}\\
\hline
\EncabezadoTabla{Campo} & \EncabezadoTabla{Contenido} \\ \hline
\endfirsthead
\multicolumn{2}{l}{\small\textit{DEC-14 — Origen de residuos sin app del contratista — continuación}}\\
\hline
\EncabezadoTabla{Campo} & \EncabezadoTabla{Contenido} \\ \hline
\endhead
\multicolumn{2}{r}{\small\textit{}}\\
\endfoot
\endlastfoot
\rowcolor{RFclaro}
Fuente de la decisión & C07, sec. 16.1, decisión 14. \\ \hline
Requerimientos relacionados & RF-AMB-01 a RF-AMB-05, RF-AUT-06 \\ \hline
\rowcolor{RFclaro}
Criterio formulado para revisión & Registrar por personal autorizado de portería/varadero o acopio la declaración del contratista, vinculada a trabajo y embarcación; identificar el recipiente y conciliar entrega al gestor. Ofrecer portal como alternativa, sin exigir instalación de app al generador. Registrar origen desconocido como discrepancia a resolver. \\ \hline
Fundamento & Se ajusta al punto de control que sí opera y permite evidencia por generador. \\ \hline
\rowcolor{RFclaro}
Impacto si el criterio resulta incorrecto & Residuos sin origen o declaración de origen sin respaldo documental. \\ \hline
Instancia y responsables propuestos de validación & Varadero y responsable ambiental: recorrer generación, depósito, traslado y certificado con un contratista sin app. \\ \hline
\rowcolor{RFclaro}
Estado de formulación & Complemento propuesto para revisión. No se presenta como acuerdo del equipo. \\ \hline
Referencia faltante & Evidencia por incorporar: SD3/SD4: captura asistida; formatos sanitarios y trazabilidad en SD5. \\ \hline
\end{longtable}
\FuenteTabla{Registro de Synaptix; fuentes y vínculos identificados en las filas.}

El impacto y la instancia de validación permiten priorizar el control de este antecedente.
\subsection*{DEC-15 — Control en portería}
La ficha establece el criterio de trabajo y las condiciones necesarias para su validación.
La Tabla~\ref{tab:a2-3-20} detalla los campos y las condiciones de validación de esta ficha.
\begin{longtable}{|L{.25\linewidth}|L{\dimexpr.75\linewidth-4\tabcolsep-3\arrayrulewidth\relax}|}
\caption{DEC-15 — Control en portería}\label{tab:a2-3-20}\\
\hline
\EncabezadoTabla{Campo} & \EncabezadoTabla{Contenido} \\ \hline
\endfirsthead
\multicolumn{2}{l}{\small\textit{DEC-15 — Control en portería — continuación}}\\
\hline
\EncabezadoTabla{Campo} & \EncabezadoTabla{Contenido} \\ \hline
\endhead
\multicolumn{2}{r}{\small\textit{}}\\
\endfoot
\endlastfoot
\rowcolor{RFclaro}
Fuente de la decisión & C07, sec. 16.1, decisión 15. \\ \hline
Requerimientos relacionados & RF-CTR-01 a RF-CTR-07 \\ \hline
\rowcolor{RFclaro}
Criterio formulado para revisión & Usar portería como control de habilitación de trabajadores y de obligaciones documentales: identidad, seguros vigentes, inducción y vínculo con faena autorizada. Dar al guardia resultado mínimo y motivo; registrar excepciones por responsable habilitado, sin sustituir procedimientos de emergencia ni bloquear rescate. \\ \hline
Fundamento & Es el punto donde la marina puede exigir antecedentes antes de una faena. \\ \hline
\rowcolor{RFclaro}
Impacto si el criterio resulta incorrecto & Ingreso no autorizado o guardia obligado a interpretar documentos complejos. \\ \hline
Instancia y responsables propuestos de validación & Portería y operaciones: prueba de vigencia, trabajador no registrado y excepción autorizada. \\ \hline
\rowcolor{RFclaro}
Estado de formulación & Criterio documentado; interfaz pendiente. No se presenta como acuerdo del equipo. \\ \hline
Referencia faltante & Evidencia por incorporar: SD3/SD4: matriz de habilitación y compatibilidad con control existente. \\ \hline
\end{longtable}
\FuenteTabla{Registro de Synaptix; fuentes y vínculos identificados en las filas.}

El impacto y la instancia de validación permiten priorizar el control de este antecedente.
\subsection*{DEC-16 — Formación ambiental externa}
La ficha establece el criterio de trabajo y las condiciones necesarias para su validación.
La Tabla~\ref{tab:a2-3-21} detalla los campos y las condiciones de validación de esta ficha.
\begin{longtable}{|L{.25\linewidth}|L{\dimexpr.75\linewidth-4\tabcolsep-3\arrayrulewidth\relax}|}
\caption{DEC-16 — Formación ambiental externa}\label{tab:a2-3-21}\\
\hline
\EncabezadoTabla{Campo} & \EncabezadoTabla{Contenido} \\ \hline
\endfirsthead
\multicolumn{2}{l}{\small\textit{DEC-16 — Formación ambiental externa — continuación}}\\
\hline
\EncabezadoTabla{Campo} & \EncabezadoTabla{Contenido} \\ \hline
\endhead
\multicolumn{2}{r}{\small\textit{}}\\
\endfoot
\endlastfoot
\rowcolor{RFclaro}
Fuente de la decisión & C07, sec. 16.1, decisión 16. \\ \hline
Requerimientos relacionados & RF-CTR-03, RF-CTR-02 \\ \hline
\rowcolor{RFclaro}
Criterio formulado para revisión & Proponer un servicio de formación por personal competente coordinado por el adjudicatario, con lista nominativa, contenidos, evaluación y vencimiento. La renovación periódica y por cambio relevante debe quedar definida con sustento; no inventar una vigencia anual. Antes de habilitar nuevas faenas se comprueba la acreditación exigible. \\ \hline
Fundamento & No existe un formador interno disponible al cual trasladar la obligación. \\ \hline
\rowcolor{RFclaro}
Impacto si el criterio resulta incorrecto & Registro de cursos sin enseñanza efectiva o costo de capacitación omitido. \\ \hline
Instancia y responsables propuestos de validación & Responsable ambiental y de servicio: revisar programa, competencia del impartidor y mecanismo de evaluación para 140 trabajadores. \\ \hline
\rowcolor{RFclaro}
Estado de formulación & Abierta periodicidad y prestación concreta. No se presenta como acuerdo del equipo. \\ \hline
Referencia faltante & Evidencia por incorporar: Plan de capacitación, contrato del impartidor, periodicidad y costos de SD11. \\ \hline
\end{longtable}
\FuenteTabla{Registro de Synaptix; fuentes y vínculos identificados en las filas.}

El impacto y la instancia de validación permiten priorizar el control de este antecedente.
\subsection*{DEC-17 — Medidas de cobranza}
La ficha establece el criterio de trabajo y las condiciones necesarias para su validación.
La Tabla~\ref{tab:a2-3-22} detalla los campos y las condiciones de validación de esta ficha.
\begin{longtable}{|L{.25\linewidth}|L{\dimexpr.75\linewidth-4\tabcolsep-3\arrayrulewidth\relax}|}
\caption{DEC-17 — Medidas de cobranza}\label{tab:a2-3-22}\\
\hline
\EncabezadoTabla{Campo} & \EncabezadoTabla{Contenido} \\ \hline
\endfirsthead
\multicolumn{2}{l}{\small\textit{DEC-17 — Medidas de cobranza — continuación}}\\
\hline
\EncabezadoTabla{Campo} & \EncabezadoTabla{Contenido} \\ \hline
\endhead
\multicolumn{2}{r}{\small\textit{}}\\
\endfoot
\endlastfoot
\rowcolor{RFclaro}
Fuente de la decisión & C07, sec. 16.1, decisión 17 (fila transcrita como «47», medidas sobre morosos). \\ \hline
Requerimientos relacionados & RF-FAC-06, RF-FAC-07 \\ \hline
\rowcolor{RFclaro}
Criterio formulado para revisión & Registrar propuestas de gestión y someter cualquier medida que afecte el acceso a decisión formal del órgano o cargo competente de la compañía, con fundamento y respaldo aplicable. No bloquear automáticamente por tramo de deuda. Conservar expediente y evidencia de ejecución o rechazo; no ejecutar cobranza judicial. \\ \hline
Fundamento & Custodia y condición de accionista requieren autoridad y fundamento explícitos. \\ \hline
\rowcolor{RFclaro}
Impacto si el criterio resulta incorrecto & Restricción indebida de acceso o medida sin respaldo documental. \\ \hline
Instancia y responsables propuestos de validación & Administración y asesor autorizado del cliente: aprobar catálogo de medidas y atribuciones antes de habilitarlo. \\ \hline
\rowcolor{RFclaro}
Estado de formulación & Parcial: autoridad y medidas admisibles no acreditadas. No se presenta como acuerdo del equipo. \\ \hline
Referencia faltante & Evidencia por incorporar: SD3: alcance; matriz de atribuciones y fundamento de cada medida. \\ \hline
\end{longtable}
\FuenteTabla{Registro de Synaptix; fuentes y vínculos identificados en las filas.}

El impacto y la instancia de validación permiten priorizar el control de este antecedente.
\subsection*{DEC-18 — Expedientes de abandonadas}
La ficha establece el criterio de trabajo y las condiciones necesarias para su validación.
La Tabla~\ref{tab:a2-3-23} detalla los campos y las condiciones de validación de esta ficha.
\begin{longtable}{|L{.25\linewidth}|L{\dimexpr.75\linewidth-4\tabcolsep-3\arrayrulewidth\relax}|}
\caption{DEC-18 — Expedientes de abandonadas}\label{tab:a2-3-23}\\
\hline
\EncabezadoTabla{Campo} & \EncabezadoTabla{Contenido} \\ \hline
\endfirsthead
\multicolumn{2}{l}{\small\textit{DEC-18 — Expedientes de abandonadas — continuación}}\\
\hline
\EncabezadoTabla{Campo} & \EncabezadoTabla{Contenido} \\ \hline
\endhead
\multicolumn{2}{r}{\small\textit{}}\\
\endfoot
\endlastfoot
\rowcolor{RFclaro}
Fuente de la decisión & C07, sec. 16.1, decisión 18. \\ \hline
Requerimientos relacionados & RF-FAC-08, RF-FAC-09, RNF-MIG-01 \\ \hline
\rowcolor{RFclaro}
Criterio formulado para revisión & Administración consolida contratos, saldos, comunicaciones, identificación del titular, custodia, fotografías y actuaciones disponibles; registra documentos faltantes como tales. Asesor autorizado define suficiencia para una actuación; conservar íntegro el legado y la cadena de cambios y exportaciones. \\ \hline
Fundamento & Una carpeta digital vacía o parcial no respalda una actuación. \\ \hline
\rowcolor{RFclaro}
Impacto si el criterio resulta incorrecto & Pérdida de prueba, titular mal identificado o acción paralizada. \\ \hline
Instancia y responsables propuestos de validación & Administración y asesor: inventario individual de 14 casos y conciliación íntegra del origen. \\ \hline
\rowcolor{RFclaro}
Estado de formulación & Criterio documentado; integridad por demostrar. No se presenta como acuerdo del equipo. \\ \hline
Referencia faltante & Evidencia por incorporar: SD5: migración y retención; lista de documentos por expediente y autorización de acceso. \\ \hline
\end{longtable}
\FuenteTabla{Registro de Synaptix; fuentes y vínculos identificados en las filas.}

El impacto y la instancia de validación permiten priorizar el control de este antecedente.
\subsection*{DEC-19 — Ocupación de boyas}
La ficha establece el criterio de trabajo y las condiciones necesarias para su validación.
La Tabla~\ref{tab:a2-3-24} detalla los campos y las condiciones de validación de esta ficha.
\begin{longtable}{|L{.25\linewidth}|L{\dimexpr.75\linewidth-4\tabcolsep-3\arrayrulewidth\relax}|}
\caption{DEC-19 — Ocupación de boyas}\label{tab:a2-3-24}\\
\hline
\EncabezadoTabla{Campo} & \EncabezadoTabla{Contenido} \\ \hline
\endfirsthead
\multicolumn{2}{l}{\small\textit{DEC-19 — Ocupación de boyas — continuación}}\\
\hline
\EncabezadoTabla{Campo} & \EncabezadoTabla{Contenido} \\ \hline
\endhead
\multicolumn{2}{r}{\small\textit{}}\\
\endfoot
\endlastfoot
\rowcolor{RFclaro}
Fuente de la decisión & C07, sec. 16.1, decisión 19. \\ \hline
Requerimientos relacionados & RF-BOY-01 a RF-BOY-05 \\ \hline
\rowcolor{RFclaro}
Criterio formulado para revisión & Registrar inspección visual identificada desde embarcación de servicio, con boya, ocupante y hora; conciliar con reservas/declaraciones al reconectar y mostrar antigüedad y discrepancias. No afirmar conocimiento en tiempo real sin infraestructura que lo acredite ni depender del piloto opcional. \\ \hline
Fundamento & Sin energía ni cobertura, la última observación no garantiza ocupación actual. \\ \hline
\rowcolor{RFclaro}
Impacto si el criterio resulta incorrecto & Asignación doble o confianza injustificada en dato antiguo. \\ \hline
Instancia y responsables propuestos de validación & Operaciones e inspectores: salida de ocho horas y conciliación con cambios ocurridos durante ausencia. \\ \hline
\rowcolor{RFclaro}
Estado de formulación & Criterio observable disponible; piloto sin justificar. No se presenta como acuerdo del equipo. \\ \hline
Referencia faltante & Evidencia por incorporar: SD3/SD4: fuente y vigencia del estado; SD11: visitas y logística. \\ \hline
\end{longtable}
\FuenteTabla{Registro de Synaptix; fuentes y vínculos identificados en las filas.}

El impacto y la instancia de validación permiten priorizar el control de este antecedente.
\subsection*{DEC-20 — Inspección del tren de fondeo}
La ficha establece el criterio de trabajo y las condiciones necesarias para su validación.
La Tabla~\ref{tab:a2-3-25} detalla los campos y las condiciones de validación de esta ficha.
\begin{longtable}{|L{.25\linewidth}|L{\dimexpr.75\linewidth-4\tabcolsep-3\arrayrulewidth\relax}|}
\caption{DEC-20 — Inspección del tren de fondeo}\label{tab:a2-3-25}\\
\hline
\EncabezadoTabla{Campo} & \EncabezadoTabla{Contenido} \\ \hline
\endfirsthead
\multicolumn{2}{l}{\small\textit{DEC-20 — Inspección del tren de fondeo — continuación}}\\
\hline
\EncabezadoTabla{Campo} & \EncabezadoTabla{Contenido} \\ \hline
\endhead
\multicolumn{2}{r}{\small\textit{}}\\
\endfoot
\endlastfoot
\rowcolor{RFclaro}
Fuente de la decisión & C07, sec. 16.1, decisión 20. \\ \hline
Requerimientos relacionados & RF-BOY-05 a RF-BOY-07, RNF-RET-06 \\ \hline
\rowcolor{RFclaro}
Criterio formulado para revisión & Registrar por componente inspección, método, hallazgo, evidencia, recomendación y reemplazo ejecutado por el mantenedor. Adoptar un programa respaldado por personal competente y especificaciones del activo; habilitar inspección extraordinaria tras eventos definidos. No fijar umbrales de desgaste ni periodicidad sin ese respaldo. \\ \hline
Fundamento & Condiciones de mar y material no pueden reducirse a una regla arbitraria del software. \\ \hline
\rowcolor{RFclaro}
Impacto si el criterio resulta incorrecto & Fondeos sin inspección adecuada o reemplazos injustificados. \\ \hline
Instancia y responsables propuestos de validación & Responsable de operaciones y especialista de fondeo: aprobar programa y criterios antes de parametrizar alertas. \\ \hline
\rowcolor{RFclaro}
Estado de formulación & Abierta la periodicidad y los disparadores técnicos. No se presenta como acuerdo del equipo. \\ \hline
Referencia faltante & Evidencia por incorporar: Programa técnico de inspección, criterios de sustitución, SD3/SD5 y costos de traslado. \\ \hline
\end{longtable}
\FuenteTabla{Registro de Synaptix; fuentes y vínculos identificados en las filas.}

El impacto y la instancia de validación permiten priorizar el control de este antecedente.
\subsection*{DEC-21 — Caída del enlace en peak}
La ficha establece el criterio de trabajo y las condiciones necesarias para su validación.
La Tabla~\ref{tab:a2-3-26} detalla los campos y las condiciones de validación de esta ficha.
\begin{longtable}{|L{.25\linewidth}|L{\dimexpr.75\linewidth-4\tabcolsep-3\arrayrulewidth\relax}|}
\caption{DEC-21 — Caída del enlace en peak}\label{tab:a2-3-26}\\
\hline
\EncabezadoTabla{Campo} & \EncabezadoTabla{Contenido} \\ \hline
\endfirsthead
\multicolumn{2}{l}{\small\textit{DEC-21 — Caída del enlace en peak — continuación}}\\
\hline
\EncabezadoTabla{Campo} & \EncabezadoTabla{Contenido} \\ \hline
\endhead
\multicolumn{2}{r}{\small\textit{}}\\
\endfoot
\endlastfoot
\rowcolor{RFclaro}
Fuente de la decisión & C07, sec. 16.1, decisión 21. \\ \hline
Requerimientos relacionados & RF-NAV-13, RF-CON-08, RF-ESC-13, RNF-CONT-01 a RNF-CONT-06 \\ \hline
\rowcolor{RFclaro}
Criterio formulado para revisión & Mantener los procesos exigidos con persistencia local 24 horas, terminales 12 horas y boyas 8 horas; conservar IDs y orden causal, resolver conflictos por regla de negocio y sincronizar en 30 minutos sin pérdidas. Publicar funciones externas indisponibles y su alternativa, sin excluir ninguna operación mínima del caso. \\ \hline
Fundamento & Los 90 zarpes del escenario no equivalen a 90 transacciones; se debe derivar y probar la carga. \\ \hline
\rowcolor{RFclaro}
Impacto si el criterio resulta incorrecto & Pérdida o duplicación de movimientos y filas de atención sin continuidad. \\ \hline
Instancia y responsables propuestos de validación & Arquitectura y operaciones: ensayo de la mañana de peak con 24 horas sin WAN y reconexión posterior. \\ \hline
\rowcolor{RFclaro}
Estado de formulación & Criterio documentado; validación pendiente. No se presenta como acuerdo del equipo. \\ \hline
Referencia faltante & Evidencia por incorporar: SD4: topología y fallas; SD5: conciliación; SD9: prueba de carga y desconexión. \\ \hline
\end{longtable}
\FuenteTabla{Registro de Synaptix; fuentes y vínculos identificados en las filas.}

El impacto y la instancia de validación permiten priorizar el control de este antecedente.
\subsection*{DEC-22 — Administración funcional}
La ficha establece el criterio de trabajo y las condiciones necesarias para su validación.
La Tabla~\ref{tab:a2-3-27} detalla los campos y las condiciones de validación de esta ficha.
\begin{longtable}{|L{.25\linewidth}|L{\dimexpr.75\linewidth-4\tabcolsep-3\arrayrulewidth\relax}|}
\caption{DEC-22 — Administración funcional}\label{tab:a2-3-27}\\
\hline
\EncabezadoTabla{Campo} & \EncabezadoTabla{Contenido} \\ \hline
\endfirsthead
\multicolumn{2}{l}{\small\textit{DEC-22 — Administración funcional — continuación}}\\
\hline
\EncabezadoTabla{Campo} & \EncabezadoTabla{Contenido} \\ \hline
\endhead
\multicolumn{2}{r}{\small\textit{}}\\
\endfoot
\endlastfoot
\rowcolor{RFclaro}
Fuente de la decisión & C07, sec. 16.1, decisión 22. \\ \hline
Requerimientos relacionados & RF-ADM-01 a RF-ADM-05, RNF-OPE-01, RNF-SOP-03 \\ \hline
\rowcolor{RFclaro}
Criterio formulado para revisión & Proponer un modelo combinado: cliente conserva aprobación de reglas y responsables de negocio; adjudicatario presta administración funcional asistida y todas las tareas especializadas. Mantener titular y suplente por función, cobertura ante vacaciones/licencias/rotación, transferencia documentada y acceso individual. No exigir un cargo TI del cliente. \\ \hline
Fundamento & Separa autoridad de negocio de capacidad técnica y evita depender de una sola persona. \\ \hline
\rowcolor{RFclaro}
Impacto si el criterio resulta incorrecto & Parámetros desactualizados, usuarios sin baja o costos omitidos durante 36 meses. \\ \hline
Instancia y responsables propuestos de validación & Jefatura del servicio y contraparte del cliente: simular ausencia del titular y revisar carga de tareas antes de cerrar oferta. \\ \hline
\rowcolor{RFclaro}
Estado de formulación & Propuesta del modelo combinado; costo y dedicación faltantes. No se presenta como acuerdo del equipo. \\ \hline
Referencia faltante & Evidencia por incorporar: SD3: responsabilidades; modelo de operación; RACI, dedicaciones, suplencias y costo de 36 meses en SD11. \\ \hline
\end{longtable}
\FuenteTabla{Registro de Synaptix; fuentes y vínculos identificados en las filas.}

El impacto y la instancia de validación permiten priorizar el control de este antecedente.
